% Unit 083 English translation of chapter6.tex, lines 2124--2312.
% Source: Methods of Algebra, Volume 2, commit
% 9a5803ff77dd3257484cb177f851a73770a59dd3 (CC BY 4.0).
% stable-unit-id: o014.aljabr2.chapter6.tate-cohomology
% Coverage: the complete Section 6.11 on Tate cohomology; ends immediately
% before Section 6.12 on profinite group cohomology at line 2313.

% segment-id: o014.li2.chapter6.tate-cohomology.h001
\section{Tate Cohomology}\label{sec:Tate-coh}
\hypertarget{o014.aljabr2.chapter6.tate-cohomology}{ }

% segment-id: o014.li2.chapter6.tate-cohomology.p001
For a finite group $G$ and a commutative ring $\Bbbk$,
Definition--Proposition~\ref{def:group-norm} introduced the element
$\nu = \sum_{g \in G} g$ of $\Bbbk[G]^G$. It lies in the center of
$\Bbbk[G]$. Thus, for every $G$-module $M$, multiplication
$x \mapsto \nu x$ gives a $G$-module homomorphism $M \to M$, again denoted
by $\nu$.

% segment-id: o014.li2.chapter6.tate-cohomology.p002
The identities $\nu g = \nu = g \nu$ imply
$\nu M \subset M^G$ and $\mathfrak{I}M \subset M^{\nu = 0}$, where
$\mathfrak{I}$ is the augmentation ideal of $\Bbbk[G]$ and $M^{\nu=0}$ is as
in Convention~\ref{con:G-mod-equalizer}. Consequently, $\nu$ induces a
$\Bbbk$-module homomorphism
% segment-id: o014.li2.chapter6.tate-cohomology.e001
\begin{equation}\label{eqn:nu-coinv-inv}
	\nu: M_G \to M^G .
\end{equation}

% segment-id: o014.li2.chapter6.tate-cohomology.p003
Since group homology (respectively, cohomology) can be regarded as a derived
version of $M_G$ (respectively, $M^G$), the canonical homomorphism
$M_G \to M^G$ suggests a bridge between homology and cohomology. Why not, for
example, consider its cokernel in some derived sense\footnote{From the
topological point of view: in the homotopical sense.}? We set this idea aside
for the moment and define the corresponding cohomology in the classical way.

% segment-id: o014.li2.chapter6.tate-cohomology.d001
\begin{definition}[Tate cohomology]\label{def:TaHm}
% segment-id: o014.li2.chapter6.tate-cohomology.x001
	\index{Tate cohomology}
% segment-id: o014.li2.chapter6.tate-cohomology.x002
	\index[sym1]{HnGMTate@$\TaHm^n(G, M)$}
	Let $G$ be a finite group and $M$ a $G$-module. For every $n \in \Z$,
	define
% segment-id: o014.li2.chapter6.tate-cohomology.g001
	\begin{align*}
		\TaHm^n(G, M) & := \Hm^n(G, M) \quad (n \geq 1), \\
		\TaHm^0(G, M) & := M^G / \nu M \\
		& = \Coker\left[ \nu: M_G \to M^G \right] \twoheadleftarrow \Hm^0(G, M), \\
		\TaHm^{-1}(G, M) & := M^{\nu=0} / \mathfrak{I} M \\
		& = \Ker\left[ \nu: M_G \to M^G \right] \hookrightarrow \Hm_0(G, M), \\
		\TaHm^{-n}(G, M) & := \Hm_{n-1}(G, M) \quad (n \geq 2).
	\end{align*}
	All of these modules are functorial in $M$.
\end{definition}

% segment-id: o014.li2.chapter6.tate-cohomology.p004
Tate cohomology is often used in Galois theory and algebraic number theory;
one example is the $\TaHm^{-1}$ appearing in
\cite[Definition 9.6.3]{Li1}. It is due to J.\ Tate.

% segment-id: o014.li2.chapter6.tate-cohomology.p005
A useful cohomology theory generally has a long exact sequence, and Tate
cohomology is no exception.

% segment-id: o014.li2.chapter6.tate-cohomology.t001
\begin{theorem}\label{prop:Tate-coh-long}
	Let $0 \to M' \to M \to M'' \to 0$ be a short exact sequence of
	$G$-modules. There is then a long exact sequence
% segment-id: o014.li2.chapter6.tate-cohomology.q001
	\[ \cdots \to \TaHm^n(G, M) \to \TaHm^n(G, M'') \xrightarrow{\delta^n} \TaHm^{n+1}(G, M') \to \TaHm^{n+1}(G, M) \to \cdots \]
	extending infinitely in both directions. The connecting homomorphisms
	$\delta^n$ are functorial with respect to morphisms of short exact
	sequences, that is, with respect to commutative diagrams with exact rows
% segment-id: o014.li2.chapter6.tate-cohomology.g002
	\[\begin{tikzcd}
		0 \arrow[r] & M' \arrow[d] \arrow[r] & M \arrow[d] \arrow[r] & M'' \arrow[d] \arrow[r] & 0 \\
		0 \arrow[r] & L' \arrow[r] & L \arrow[r] & L'' \arrow[r] & 0
	\end{tikzcd}\]
\end{theorem}

% segment-id: o014.li2.chapter6.tate-cohomology.pf001
\begin{proof}
	Splice the long exact sequences in homology and cohomology as follows:
% segment-id: o014.li2.chapter6.tate-cohomology.g003
	\[\begin{tikzcd}[column sep=small, every cell/.append style={font = \footnotesize}]
		\cdots \arrow[r] & \Hm_1(G, M'') \arrow[r] \arrow[d] & \Hm_0(G, M') \arrow[r] \arrow[d, "{N_1}"] & \Hm_0(G, M) \arrow[r] \arrow[d, "N_2"] & \Hm_0(G, M'') \arrow[r] \arrow[d, "N_3"] & 0 \arrow[d] & \\
		& 0 \arrow[r] & \Hm^0(G, M') \arrow[r] & \Hm^0(G, M) \arrow[r] & \Hm^0(G, M'') \arrow[r] & \Hm^1(G, M') \arrow[r] & \cdots
	\end{tikzcd}\]
	Here $N_1,N_2,N_3$ are all the homomorphisms previously denoted by
	$\nu$. The diagram above is commutative and its rows are exact. The only
	point requiring explanation is the commutativity of the two squares that
	touch $0$. This can be checked explicitly using the standard complex and
	the chain complex, while recalling the connecting morphisms
	$\Hm_1 \to \Hm_0$ and $\Hm^0 \to \Hm^1$; the verification is left as an
	exercise.

	The Snake Lemma now gives the exact sequence
% segment-id: o014.li2.chapter6.tate-cohomology.g004
	\[\begin{tikzcd}[row sep=small, column sep=small, every cell/.append style={font = \footnotesize}]
		\Ker(N_1) \arrow[r] & \Ker(N_2) \arrow[r] & \Ker(N_3) \arrow[r, "{\delta^{-1}}"] & \Coker(N_1) \arrow[r] & \Coker(N_2) \arrow[r] & \Coker(N_3) \\
		\TaHm^{-1}(G, M') \arrow[equal, u] & \TaHm^{-1}(G, M) \arrow[equal, u] & \TaHm^{-1}(G, M'') \arrow[equal, u] & \TaHm^0(G, M') \arrow[equal, u] & \TaHm^0(G, M) \arrow[equal, u] & \TaHm^0(G, M'') \arrow[equal, u]
	\end{tikzcd}\]
	This is one part of the desired exact sequence, and the functoriality of
	the connecting homomorphism $\delta^{-1}$ follows from the Snake Lemma.
	The remaining parts reduce to the long exact sequences in group homology
	and group cohomology, together with the commutative diagram with exact rows
	proved in the first step.
\end{proof}

% segment-id: o014.li2.chapter6.tate-cohomology.p006
We next state the Tate-cohomology version of Shapiro's Lemma
(Theorem~\ref{prop:Shapiro}). Recall that if $G$ is finite, then for every
subgroup $H$ the two induction functors
$\Ind^G_H,\iInd^G_H:H\dcate{Mod}\to G\dcate{Mod}$ coincide; see
Corollary~\ref{prop:iInd-Ind}. From now on both will be denoted by
$\Ind^G_H$.

% segment-id: o014.li2.chapter6.tate-cohomology.pr001
\begin{proposition}
% segment-id: o014.li2.chapter6.tate-cohomology.x003
	\index{Shapiro's lemma}
	For a finite group $G$, a subgroup $H$, and an $H$-module $N$, there is a
	canonical isomorphism
% segment-id: o014.li2.chapter6.tate-cohomology.q002
	\[ \TaHm^n(H, N) \simeq \TaHm^n\left( G, \Ind^G_H(N)\right), \quad n \in \Z. \]
\end{proposition}

% segment-id: o014.li2.chapter6.tate-cohomology.pf002
\begin{proof}
	For $n \geq 1$ or $n \leq -2$, the assertion reduces to
	Theorem~\ref{prop:Shapiro}. We next handle the cases $n=-1,0$.
	Realize $\Ind^G_H(N)$ as a mapping space by
	Lemma~\ref{prop:Ind-as-mappings}. Write $\nu^G$ (respectively, $\nu^H$)
	for the $\nu$ corresponding to $G$ (respectively, $H$). It remains to
	prove that the following diagram commutes:
% segment-id: o014.li2.chapter6.tate-cohomology.g005
	\[\begin{tikzcd}
		\Ind^G_H(N)_G \arrow[r, "{\nu^G}"] & \Ind^G_H(N)^G \arrow[d, "\sim" sloped] \\
		N_H \arrow[u, "\sim" sloped] \arrow[r, "{\nu^H}"'] & N^H
	\end{tikzcd}\]
	The vertical isomorphisms come from
	Corollary~\ref{prop:Ind-invariant}. Let $y\in N$, and write $[y]$ for its
	image in $N_H$. The image of $[y]$ in $\Ind^G_H(N)_G$ is $[f]$, where
	$f\in\Ind^G_H(N)$ satisfies $f(1_G)=y$ and
	$\Supp(f)\subset H$. Thus the image in $N^H$ of
	$\nu^Gf:x\mapsto\sum_{g\in G}f(xg)$ is
	$(\nu^Gf)(1_G)=\sum_gf(g)$. But this also equals
	$\sum_{h\in H}f(h)=\sum_hhf(1_G)=\nu^Hy$. This proves the assertion.
\end{proof}

% segment-id: o014.li2.chapter6.tate-cohomology.p007
The argument can be strengthened to show that the canonical isomorphism is
compatible with the long exact sequences on both sides.

% segment-id: o014.li2.chapter6.tate-cohomology.c001
\begin{corollary}\label{prop:Tate-ind-vanishing}
	Let $G$ be a finite group. For every $\Bbbk$-module $N$,
% segment-id: o014.li2.chapter6.tate-cohomology.q003
	\[ \TaHm^n\left(G, \Ind^G_{\{1\}}(N)\right) = 0, \quad n \in \Z. \]
\end{corollary}

% segment-id: o014.li2.chapter6.tate-cohomology.pf003
\begin{proof}
	Clearly $\TaHm^n(\{1\},N)=0$ for every $n\in\Z$.
\end{proof}

% segment-id: o014.li2.chapter6.tate-cohomology.ex001
\begin{example}\label{eg:additive-90-finite}
% segment-id: o014.li2.chapter6.tate-cohomology.x004
	\index{Hilbert's Theorem 90}
	Let $E|F$ be a finite Galois extension of fields, with Galois group
	$\Gal(E|F)$ acting on $E$. Take the coefficient ring to be
	$\Bbbk=F$ or $\Z$. The Normal Basis Theorem
	\cite[Theorem 9.5.6]{Li1} is equivalent to the statement
	$E\simeq\Ind^{\Gal(E|F)}_{\{1\}}(F)$, so
	Corollary~\ref{prop:Tate-ind-vanishing} implies
% segment-id: o014.li2.chapter6.tate-cohomology.q004
	\[ \TaHm^n(\Gal(E|F), E) = 0, \quad n \in \Z; \]
	therefore $n\geq1\implies\Hm^n(\Gal(E|F),E)=0$. If $E|F$ is cyclic and
	$n=-1$, this gives another proof of the additive version of Hilbert's
	Theorem 90 mentioned in \cite[Theorem 9.6.10]{Li1}. For the
	multiplicative version, see Example~\ref{eg:multiplicative-90-finite}.
\end{example}

% segment-id: o014.li2.chapter6.tate-cohomology.p008
For a general $H\subset G$ and a $G$-module $M$, the left and right adjoints
of restriction give canonical morphisms
% segment-id: o014.li2.chapter6.tate-cohomology.q005
\[ \epsilon_M: \iInd^G_H(\Res^G_H(M)) \to M, \quad \eta_M: M \to \Ind^G_H(\Res^G_H(M)). \]

% segment-id: o014.li2.chapter6.tate-cohomology.p009
From the concrete description of the adjoint pairs
(Proposition~\ref{prop:pullback-two-adjoint}),
$\epsilon_M(g\otimes x)=gx$, while $\eta_M(x)=[g\mapsto gx]$.
Thus $\epsilon_M$ is surjective and $\eta_M$ is injective. Returning to the
case of finite groups, we have
$\iInd^G_{\{1\}}=\Ind^G_{\{1\}}$. Using $\epsilon_M$ and $\eta_M$, the
vanishing in Corollary~\ref{prop:Tate-ind-vanishing} shows that each
$\TaHm^n(G,\cdot)$ is both erasable and coerasable
(Definition~\ref{def:erasable}). Consequently, the dimension-shifting
technique of Proposition~\ref{prop:dimension-shifting} has the following
simplified form.

% segment-id: o014.li2.chapter6.tate-cohomology.c002
\begin{corollary}[Dimension shifting]\label{prop:Tate-dim-shift}
	For every $G$-module $M$, as usual write $M$ as an abbreviation for
	$\Res^G_{\{1\}}M$, and consider the short exact sequences
% segment-id: o014.li2.chapter6.tate-cohomology.g006
	\begin{gather*}
		0 \to M' \to \Ind^G_{\{1\}}(M) \xrightarrow{\epsilon_M} M \to 0, \\
		0 \to M \xrightarrow{\eta_M} \Ind^G_{\{1\}}(M) \to M'' \to 0.
	\end{gather*}
	The corresponding connecting homomorphisms in the long exact sequences
	give canonical isomorphisms
% segment-id: o014.li2.chapter6.tate-cohomology.q006
	\[ \TaHm^{n-1}(G, M'') \rightiso \TaHm^n(G, M) \rightiso \TaHm^{n+1}(G, M'), \quad n \in \Z. \]
\end{corollary}

% segment-id: o014.li2.chapter6.tate-cohomology.pf004
\begin{proof}
	Combine the long exact sequence
	(Theorem~\ref{prop:Tate-coh-long}) with the vanishing in
	Corollary~\ref{prop:Tate-ind-vanishing}.
\end{proof}

% segment-id: o014.li2.chapter6.tate-cohomology.p010
Many properties of Tate cohomology can be reduced by dimension shifting to
the case $n=0$.

% segment-id: o014.li2.chapter6.tate-cohomology.p011
Return to the definition. Choose an arbitrary projective resolution $P\to M$
and injective resolution $M\to I$ of the $G$-module $M$, regard both as
complexes, and then consider the morphism of complexes $P_G\to I^G$, shown
in the diagram
% segment-id: o014.li2.chapter6.tate-cohomology.g007
\[\begin{tikzcd}
	\cdots \arrow[r] & 0 \arrow[r] & 0 \arrow[r] & I^{0, G} \arrow[r] & I^{1, G} \arrow[r] & I^{2, G} \arrow[r] & \cdots \\
	\cdots \arrow[r] & P_{2, G} \arrow[r] \arrow[u] & P_{1, G} \arrow[r] \arrow[u] & P_{0, G} \arrow[u, "\Omega"] \arrow[r] & 0 \arrow[r] \arrow[u] & 0 \arrow[r] \arrow[u] & \cdots
\end{tikzcd}\]
Here $\Omega$ is the composite
$P_{0,G}\twoheadrightarrow M_G\xrightarrow{\nu}M^G\hookrightarrow I^{0,G}$;
commutativity is clear. Take the mapping cone
$\Cone\left[P_G\to I^G\right]$, or equivalently the total complex of the
diagram above (placing $P_G$ in vertical coordinate $-1$). The reader is
asked to verify that
% segment-id: o014.li2.chapter6.tate-cohomology.e002
\begin{equation}\label{eqn:Tate-cone-1}
	\Hm^n\left( \Cone\left[P_G \to I^G\right] \right) = \TaHm^n(G, M).
\end{equation}

% segment-id: o014.li2.chapter6.tate-cohomology.p012
The equation above in fact shows how to define a canonical morphism
$\Bbbk \otimesL[{\Bbbk[G]}] M \to \RHom_G(\Bbbk, M)$ in the derived category
$\cate{D}(G)$ of $G$-modules. Complete this morphism to a distinguished
triangle in $\cate{D}(G)$,
% segment-id: o014.li2.chapter6.tate-cohomology.q007
\[ \Bbbk \otimesL[{\Bbbk[G]}] M \to \RHom_G(\Bbbk, M) \to \mathrm{Tate}(G, M) \xrightarrow{+1}, \]
so that $\Hm^n(\mathrm{Tate}(G,M))\simeq\TaHm^n(G,M)$. The difficulty with
this construction is that different choices for $\mathrm{Tate}(G,M)$ in the
distinguished triangle are isomorphic, but the isomorphisms are not
canonical. This is a conspicuous defect of the theory of derived categories;
see Corollary~\ref{prop:Cone-uniqueness} and the explanation following it.

% segment-id: o014.li2.chapter6.tate-cohomology.p013
Nevertheless, suppose that
$\Bbbk \otimesL[{\Bbbk[G]}] M$ (respectively, $\RHom_G(\Bbbk,M)$) is
represented by a chosen complex concentrated in nonpositive (respectively,
nonnegative) degrees, including but not limited to $P_G$ (respectively, $I^G$)
in \eqref{eqn:Tate-cone-1}. We can still define the morphism $\Omega$, take
its mapping cone as $\mathrm{Tate}(G,M)$, and interpret it as the homotopy
cokernel in the sense of Remark~\ref{rem:homotopy-kernel-cokernel}.\footnote{Taking
the homotopy kernel of
$\Bbbk \otimesL[{\Bbbk[G]}] M \to \RHom_G(\Bbbk, M)$ produces no new
structure; it differs from the homotopy cokernel only by a shift.} The
exercises will explain how to use the standard resolution of $\Bbbk$ instead,
providing further tools for computing Tate cohomology.

% segment-id: o014.li2.chapter6.tate-cohomology.p014
We now discuss the special case of finite cyclic groups in greater detail.

% segment-id: o014.li2.chapter6.tate-cohomology.ex002
\begin{example}\label{eg:TaHm-periodic}
	Let $m\in\Z_{\geq1}$ and let $C_m$ be the cyclic group of order $m$;
	choose any generator $\sigma$ of $C_m$. For every $C_m$-module $A$, we
	claim that $\TaHm^n(C_m,A)$ is the $\Hm^n$ of the following complex
	($n\in\Z$):
% segment-id: o014.li2.chapter6.tate-cohomology.q008
	\[ \cdots \xrightarrow{\sigma-1} A \xrightarrow{\nu} \underbracket{A}_{\text{degree-zero term}} \xrightarrow{\sigma-1} A \xrightarrow{\nu} A \xrightarrow{\sigma-1} \to \cdots \]
	This complex repeats with period $2$. Indeed, the proof of
	Proposition~\ref{prop:group-cohomology-cyclic} explicitly wrote down the
	complexes computing homology and cohomology; splicing them by $\nu$ yields
	the mapping cone $\mathrm{Tate}(G,M)$ mentioned above, in precisely the
	form displayed. We immediately obtain canonical isomorphisms
% segment-id: o014.li2.chapter6.tate-cohomology.q009
	\[ \TaHm^n(C_m, A) \simeq \begin{cases}
		A^{\nu=0} / (\sigma-1) A, & n\;\text{odd} \\
		A^{\sigma=1} / \nu A, & n\;\text{even}.
	\end{cases}\]

	Periodicity shows that if
	$0\to A'\to A\to A''\to0$ is a short exact sequence of $C_m$-modules,
	then the long exact sequence of Theorem~\ref{prop:Tate-coh-long} fits into
	the exact hexagon
% segment-id: o014.li2.chapter6.tate-cohomology.e003
	\begin{equation}\label{eqn:Herbrand-hexagon}
		\begin{tikzcd}[column sep=small]
			& \TaHm^{\text{even}}(C_m, A') \arrow[r, "f_1"] & \TaHm^{\text{even}}(C_m, A) \arrow[rd, "f_2"] & \\
			\Hm^{\text{odd}}(C_m, A'') \arrow[ru, "f_6"] & & & \TaHm^{\text{even}}(C_m, A'') \arrow[ld, "f_3"] . \\
			& \TaHm^{\text{odd}}(C_m, A) \arrow[lu, "f_5"] & \TaHm^{\text{odd}}(C_m, A') \arrow[l, "f_4"] &
		\end{tikzcd}
	\end{equation}
\end{example}

% segment-id: o014.li2.chapter6.tate-cohomology.d002
\begin{definition}
% segment-id: o014.li2.chapter6.tate-cohomology.x005
	\index{Herbrand quotient}
	Let $A$ be a $C_m$-module ($m\in\Z_{\geq1}$). Assuming that
	$\TaHm^n(C_m,A)$ is finite for every $n$, define the positive rational
	number
% Editorial correction carried from the Indonesian witness: the Chinese source
% says "positive integer," but the displayed quotient need only be rational.
% segment-id: o014.li2.chapter6.tate-cohomology.q010
	\[ h(C_m, A) := \frac{\left|\TaHm^{\text{even}}(C_m, A)\right|}{\left|\TaHm^{\text{odd}}(C_m, A)\right|}, \]
	called the \emph{Herbrand quotient} of $A$.
\end{definition}

% segment-id: o014.li2.chapter6.tate-cohomology.p015
If $A$ is finite, then $\TaHm^n(C_m,A)$ is clearly finite. Note that the
number of its elements is independent of the choice of coefficient ring
$\Bbbk$; we may take $\Bbbk=\Z$. In algebraic number theory, the following
theorem of Herbrand is a fundamental and important result.

% segment-id: o014.li2.chapter6.tate-cohomology.t002
\begin{theorem}[J.\ Herbrand]
	Let $m\in\Z_{\geq1}$ and let $C_m$ be the cyclic group of order $m$.
% segment-id: o014.li2.chapter6.tate-cohomology.l001
	\begin{enumerate}[(i)]
% segment-id: o014.li2.chapter6.tate-cohomology.i001
		\item Let $0\to A'\to A\to A''\to0$ be a short exact sequence of
		$C_m$-modules. If two of $h(C_m,A')$, $h(C_m,A)$, and
		$h(C_m,A'')$ are defined, then the third is automatically defined; in
		that case,
% segment-id: o014.li2.chapter6.tate-cohomology.q011
		\[ h(C_m, A) = h(C_m, A') h(C_m, A''). \]
% segment-id: o014.li2.chapter6.tate-cohomology.i002
		\item If $A$ is a finite $C_m$-module, then $h(C_m,A)=1$.
	\end{enumerate}
\end{theorem}

% segment-id: o014.li2.chapter6.tate-cohomology.pf005
\begin{proof}
	For~(i), consider the exact hexagon
	\eqref{eqn:Herbrand-hexagon}. Write
	$n_i:=|\Image(f_i)|$. In the sense of multiplication of cardinalities, the
	cardinalities of the six terms are, respectively,
% segment-id: o014.li2.chapter6.tate-cohomology.g008
	\begin{center}\begin{tikzpicture}[commutative diagrams/every diagram,
		one/.style={draw, rectangle, minimum size=6mm, rounded corners=3mm},
		two/.style={draw, rectangle, minimum size=6mm}
		]
		\node[one] (P0) at (60:1.75cm) {$n_1 n_2$};
		\node[two] (P1) at (60+60:1.75cm) {$n_6 n_1$} ;
		\node[one] (P2) at (60+2*60:1.75cm) {$n_5 n_6$};
		\node[two] (P3) at (60+3*60:1.75cm) {$n_4 n_5$};
		\node[one] (P4) at (60+4*60:1.75cm) {$n_3 n_4$};
		\node[two] (P5) at (0:1.75cm) {$n_2 n_3$};
	\end{tikzpicture}\end{center}

	The three Herbrand quotients correspond to the quotients of opposite
	terms in the diagram above. Thus assertion~(i) is equivalent to the
	following observations:
% segment-id: o014.li2.chapter6.tate-cohomology.l002
	\begin{itemize}
% segment-id: o014.li2.chapter6.tate-cohomology.i003
		\item if the cardinalities at both ends of two diagonals are finite,
		then the same holds at both ends of the remaining diagonal;
% segment-id: o014.li2.chapter6.tate-cohomology.i004
		\item the product of the terms in square-cornered boxes equals the
		product of the terms in round-cornered boxes.
	\end{itemize}

	For~(ii), there are short exact sequences of finite $\Z$-modules
% segment-id: o014.li2.chapter6.tate-cohomology.q012
	\[ 0 \to A^{\sigma=1} \to A \xrightarrow{\sigma-1} \mathfrak{I} A \to 0, \quad 0 \to A^{\nu=0} \to A \xrightarrow{\nu} \nu A \to 0; \]
	here $(\sigma-1)A=\mathfrak{I}A$ follows from
	Lemma~\ref{prop:cyclic-resolution-prep}. Therefore,
% segment-id: o014.li2.chapter6.tate-cohomology.q013
	\[ |A^{\sigma=1}| \cdot |\mathfrak{I} A| = |A^{\nu=0}| \cdot |\nu A|, \]
	which rearranges to
	$|\TaHm^0(C_m,A)|=|\TaHm^{-1}(C_m,A)|$. This proves the assertion.
\end{proof}
